\section{Institutional Circulation and Workshop Guide}
\label{sec:circulation-guide}

This section supports responsible circulation to schools, DPOs, ethics committees, educational
leadership, and Horizon Europe partners.
It is not a marketing summary; it describes \textbf{how} to discuss the pilot so communities can
make informed, reversible decisions.

\subsection{Recommended document routing}

\begin{table}[htbp]
  \centering
  \caption{Which document to share with whom (first contact).}
  \label{tab:document-routing}
  \small
  \begin{tabular}{@{}p{2.8cm}p{2.4cm}p{5.8cm}@{}}
    \toprule
    Audience & Start with & Then, if proceeding \\
    \midrule
    Leadership / board & Executive Brief & Charter workshop; full requirements pack \\
    DPO / legal & Executive Brief + Box~1 & Full DPIA checklist (Section~\ref{sec:dpia}) \\
    Ethics committee & Executive Brief & Transparency plan; evaluation instruments \\
    Safety coordinators & Requirements body & SOPs; protocol matrix (Appendix~\ref{app:governance-templates}) \\
    Parent / staff reps & Transparency notice outline & Q\&A session; no technical appendix by default \\
    Horizon partners & Executive Brief & Governance harmonization; anonymized reporting rules \\
    IT (when engaged) & Trust-zone summary & Technical reference (Appendix~\ref{app:technical-architecture}) \\
    \bottomrule
  \end{tabular}
\end{table}

\subsection{Workshop sequence (recommended)}

\begin{enumerate}[noitemsep]
  \item \textbf{Steering formation} (leadership, safety, DPO): scope, boundaries (Box~1), pause/stop rights.
  \item \textbf{Protocol-mapping workshop} (safety, admin, facilities): map zone summary types to
        existing SMS branches \textbf{before} any live sensing (Appendix~\ref{app:governance-templates}).
  \item \textbf{Community transparency session} (as required locally): plain-language notice; Q\&A;
        document concerns and scope adjustments.
  \item \textbf{Sandbox familiarization} (safety + IT): synthetic feeds only; advisory-only practice.
  \item \textbf{Tabletop exercises} ($\geq$2): role-play with sandboxed inputs; record remediations.
  \item \textbf{Go/no-go review} (steering + DPO): checklist in Section~\ref{sec:pilot-programme}.
\end{enumerate}

\subsection{Phased pilot communication (plain language)}

\begin{table}[htbp]
  \centering
  \caption{Communication milestones by phase.}
  \label{tab:phased-comms}
  \small
  \begin{tabular}{@{}p{1.2cm}p{2.6cm}p{5.2cm}@{}}
    \toprule
    Phase & Audience & Message (honest framing) \\
    \midrule
    0 & Internal & ``We are \emph{planning} a governed pilot; no live sensing yet.'' \\
    1 & Staff / parents (if required) & ``Training and tabletop only; fictional or sandboxed inputs.'' \\
    2 & Community & ``Limited zones; staffed-hours review; you may ask to pause or exclude a zone.'' \\
    3+ & Consortium & ``Anonymized learning only; schools retain operational control.'' \\
    Any & All & ``This supports staff review; it does not replace your emergency plan.'' \\
    \bottomrule
  \end{tabular}
\end{table}

Avoid promising harm reduction, compliance certification, or permanent installation at Phase~0--1.

\subsection{Building community trust (operational practices)}

\begin{itemize}[noitemsep]
  \item \textbf{Co-design by default:} schools choose zones with safety staff; sensitive spaces need
        explicit steering approval.
  \item \textbf{Transparency before expansion:} update notices and signage when zones or modalities change.
  \item \textbf{Parent and staff voice:} consult councils or unions per local law \emph{before} Phase~2.
  \item \textbf{Right to pause:} document how any stakeholder can trigger review or stop conditions.
  \item \textbf{Local processing:} structured summaries remain on school-controlled systems unless
        governance approves otherwise (Section~\ref{sec:deployment-infrastructure}).
  \item \textbf{No surprise capabilities:} biometrics, identity recognition, or autonomous actuation
        are excluded by design (Box~\ref{box:system-boundaries})---not future ``feature toggles.''
\end{itemize}

\subsection{Likely stakeholder concerns (prepare responses)}

\begin{table}[htbp]
  \centering
  \caption{Anticipated objections and governance-aligned responses.}
  \label{tab:stakeholder-objections}
  \small
  \begin{tabular}{@{}p{3.2cm}p{5.8cm}@{}}
    \toprule
    Concern & Response posture \\
    \midrule
    ``This is surveillance.'' &
      Zone-level summaries only; no identity outputs; co-design; opt-out; DPO oversight. \\
    ``Cameras will stay forever.'' &
      Phased pilot with end date; steering decides extend/modify/halt. \\
    ``The system will lock down the school.'' &
      Advisory-only; staff retain authority; tabletop proves handoffs. \\
    ``Children's privacy.'' &
      DPIA; minimization; transparency; DPO sign-off before live sensing. \\
    ``Research partners control operations.'' &
      Research supports training/evaluation; school owns disposition and SMS. \\
    ``False alarms will overwhelm staff.'' &
      Staffed-hours scope; calibration; stop thresholds; bounded review queues. \\
    \bottomrule
  \end{tabular}
\end{table}

\subsection{Governance gaps institutions should close locally}

The pack provides planning structure; each school must still assign:
\begin{itemize}[noitemsep]
  \item named \textbf{steering chair} and \textbf{safety lead} with backup/substitute;
  \item \textbf{lawful basis} and children's safeguards in a signed DPIA;
  \item \textbf{union / works council} consultation where mandatory;
  \item \textbf{insurance and vendor DPAs} for any processors;
  \item \textbf{incident playbooks} linking stop conditions to communications (parents, staff, regulators);
  \item \textbf{evaluation ethics approval} if staff interviews or surveys are used.
\end{itemize}

Templates for charter, transparency notice, protocol matrix, and tabletop record:
Appendix~\ref{app:governance-templates}.
